\chapter*{Acronymes}
\addcontentsline{toc}{chapter}{Acronymes}

\begin{acronym}[ASVS-XXX]
\acro{AD}{Active Directory --- annuaire d'identités Microsoft}
\acro{ANSSI}{Agence nationale de la sécurité des systèmes d'information}
\acro{API}{Application Programming Interface}
\acro{ASVS}{Application Security Verification Standard (OWASP)}
\acro{AQSSI}{Autorité qualifiée pour la sécurité des SI}
\acro{AZ}{Azure CLI (outil ligne de commande Microsoft Azure)}
\acro{CI}{Continuous Integration --- intégration continue}
\acro{CLI}{Command-Line Interface}
\acro{CVE}{Common Vulnerabilities and Exposures}
\acro{CVSS}{Common Vulnerability Scoring System}
\acro{DC}{Domain Controller --- contrôleur de domaine Active Directory}
\acro{EBIOS RM}{EBIOS Risk Manager --- méthode ANSSI d'analyse de risque}
\acro{ENT}{Espace Numérique de Travail (universités)}
\acro{HOMO-CI}{Homologation Continuous Integration --- nom de l'outil}
\acro{IDOR}{Insecure Direct Object Reference}
\acro{IOC}{Indicator of Compromise}
\acro{ISO 27001}{Norme internationale de management de la sécurité de l'information}
\acro{JSON}{JavaScript Object Notation}
\acro{JSONL}{JSON Lines --- un objet JSON par ligne}
\acro{LDAP}{Lightweight Directory Access Protocol}
\acro{MFA}{Multi-Factor Authentication}
\acro{NSG}{Network Security Group (Azure)}
\acro{OWASP}{Open Worldwide Application Security Project}
\acro{PoC}{Proof of Concept --- preuve de concept}
\acro{RGPD}{Règlement général sur la protection des données}
\acro{RGS}{Référentiel Général de Sécurité (France)}
\acro{SHA-256}{Secure Hash Algorithm 256 bits}
\acro{SI}{Système d'Information}
\acro{SIEM}{Security Information and Event Management}
\acro{SOC}{Security Operations Center}
\acro{SQL}{Structured Query Language}
\acro{TLS}{Transport Layer Security}
\acro{UUID}{Universally Unique Identifier}
\acro{VM}{Virtual Machine --- machine virtuelle}
\acro{XSS}{Cross-Site Scripting}
\end{acronym}
